\chapter{AgentOS 双轴多尺度总体架构}
\label{chap:agentos-dual-axis-multiscale}

\begin{quote}
\textbf{【本章总体说明】}

本章对前述三相记忆、工作记忆有限性、SPC--AHC--TCE 认知调节、CCM、Scheduler、Boundary、认知卸载、Body Runtime、SGC/FCS、BPCC、技能编译与 KRA 等机制进行一次结构性的统一。我们不再把 AgentOS 理解为若干模块的并列集合，而把它严格描述为一个由\textbf{功能层级轴（Functional Hierarchy）}与\textbf{时间尺度轴（Temporal Hierarchy）}共同张成的多尺度认知动力系统。任何 Agent-CSO 在系统中的运行状态，都必须同时回答两个问题：第一，它当前由哪个功能结构承载；第二，它当前以什么时间尺度存在和变化。由此，本章确立 AgentOS 的双轴基本原理：
\[
\boxed{
AgentOS
=
FunctionalHierarchy
\times
TemporalHierarchy
}
\]
并进一步说明，持续智能的学习过程并非简单的参数增量，而主要表现为 Agent-CSO 在该双轴空间中的持续重新配置：结构可以向更慢时间尺度沉降，也可以由慢结构重新激活进入快速过程；异常和新颖结构可以从局部功能向更高层级提升，而成熟能力则可以从显式认知向更低、更局部、更快速的执行闭环编译。最终，AgentOS 的运行可以概括为四种基本跨尺度流：\textbf{向慢巩固、向快激活、向上升级、向下编译}。

本章所讨论的是一个固定总体功能骨架内部的多尺度动力学原理，而非任意自修改的软件拓扑。功能边、时间常数、承载位置和局部策略可以发生受控变化，但智能主体、现实验证路径以及安全边界保持结构连续。因而，本章既是前述各模块从局部机制走向统一系统动力学的收束，也是理解 AgentOS 为什么能够在有限工作记忆条件下形成长期学习、技能成熟和选择性显式认知的关键章节。
\end{quote}


\section{功能层级：智能不是单一中央处理器}

如果我们只从传统软件工程的角度观察 AgentOS，很容易把系统画成一组模块：工作记忆 $h$、情景记忆 $E$、结构记忆 $\Theta$、自我 $\Psi$、SPC、AHC、TCE、CCM、Scheduler、Body Runtime、BPCC，以及不同外部工具。这样的组件图对于工程实现是必要的，却不足以解释智能为何能够持续运行。真正需要理解的是：这些组件并非平等地、同步地围绕一个中央推理器工作，而是在不同功能层级中关闭各自的局部反馈环，并仅在局部环无法处理当前结构时才发生跨层传播。

因此，我们首先把 AgentOS 的功能系统定义为一个带有多重闭环的有向功能图
\[
\mathcal G_F
=
(V_F,E_F,\mathcal L),
\]
其中 $V_F=\{F_1,\ldots,F_n\}$ 表示功能节点，$E_F$ 表示功能之间可以发生状态、控制和 Agent-CSO 迁移的作用关系，而 $\mathcal L=\{L_1,\ldots,L_m\}$ 表示系统内部实际能够独立闭合的反馈回路。这里最重要的一点是，一个节点并不必然只属于一个闭环。SPC 可以同时参与工作记忆运行控制、自我感知和心理稳定性闭环；Body Runtime 可以同时参与感知、控制、技能和安全闭环。因此更加一般地可以写成
\[
F_i\in L_a\cap L_b\cap L_c.
\]
这意味着 AgentOS 更接近一个闭环超图，而不是一棵严格的树。

尽管如此，为了理解跨尺度控制，我们仍然可以定义一个功能偏序
\[
F_i\prec_F F_j,
\]
表示在某一类控制关系中，$F_i$ 更接近局部、快速和具体的执行，而 $F_j$ 更接近全局、抽象和长期约束。这里的“高层”和“低层”并不是价值判断，也不是软件调用栈意义上的上层 API 与下层 API，而是表示\textbf{一个功能作用所覆盖的状态空间大小以及它需要协调的自由度数量}。

例如，BPCC 对具体执行器轨迹进行局部预测控制，其自由度范围相对有限；工作记忆 $h$ 同时容纳当前目标、重要对象、局部计划和跨模块约束，其协调范围更大；$\Theta$ 则保存能够在大量未来场景中重复使用的生成结构；$\Psi$ 进一步对什么结构值得保留、什么经历属于“我”、什么目标与身份一致施加超慢约束。于是可以形成近似的功能层级：
\[
BodyFastLoop
\prec_F
h
\prec_F
\Theta
\prec_F
\Psi.
\]

需要强调的是，这不是一种强制所有信息逐层经过的流水线。成熟 AgentOS 的正常状态恰恰应该避免这种中央化。例如，机器人已经学会保持平衡以后，其身体姿态扰动通常应在 BPCC 或更低级控制器中直接解决，而不应该形成完整 Agent-CSO 后再进入 $h$ 进行显式推理。只有当局部闭环的预测残差持续增大，或者当前异常与高层目标产生关系时，相关结构才有必要向上升级。因此：

\[
\boxed{
EachScaleShouldCloseItsOwnFastLoop
}
\]

成为整个功能层级的基本原则。

从这一点出发，我们可以重新理解所谓“智能成熟”。幼稚系统往往具有高度中央化的运行方式：大量细节必须经过同一个模型、同一个上下文和同一个推理循环；成熟系统则逐渐形成大量局部自治闭环，只把少数真正无法局部闭合的问题提升到显式认知层。因而：

\[
\boxed{
MatureIntelligence
\neq
MoreCentralReasoning
}
\]

而更加接近：

\[
\boxed{
MatureIntelligence
=
MoreLocalClosure
+
MoreSelectiveGlobalCognition.
}
\]

这正是功能层级存在的根本意义。它不是为了增加系统结构的复杂程度，而是为了使不同类型的 Agent-CSO 能够在最适合自己的功能区域完成处理，使有限的全局认知资源只承担真正具有跨域、跨目标和跨时间尺度意义的问题。


\section{时间层级：同一智能体内部存在多种认知速度}

功能层级只能告诉我们“谁在处理”，却不能回答“以多快的速度处理”。因此 AgentOS 的第二个基本坐标是时间尺度。一个持续智能体内部并不存在单一统一的认知时钟。身体姿态控制可能在毫秒到百毫秒尺度变化，感知语义状态可以在数百毫秒到秒级变化，工作记忆中的目标和推理结构可能维持数秒到数分钟，情景经验可以在小时或天级重新组织，$\Theta$ 的稳定生成结构可能需要更长的经验才能修改，而 $\Psi$ 与 $\Omega$ 所代表的身份和生命周期方向则必须具有远慢于普通任务状态的变化速度。

因此可以为任何 Agent-CSO $C$ 定义一个有效时间尺度
\[
\tau(C),
\]
它不表示该结构“多久才执行一次”，而表示该结构在何种时间窗口内可以被视为相对稳定。若
\[
\tau(C_1)\ll\tau(C_2),
\]
则在 $C_1$ 已经经历大量变化的时间内，$C_2$ 仍然可以近似看作慢变量。

这一关系可以用前述多时间尺度动力系统表示。设快速状态为 $x_f$，慢变量为 $z_s$，则
\[
\dot x_f
=
F(x_f,z_s,u),
\]
而
\[
\dot z_s
=
\epsilon G(z_s,\langle \varepsilon_f\rangle_T),
\qquad
0<\epsilon\ll1.
\]
在快速时间尺度中，$z_s$ 近似固定，因此快速动力学是在慢结构所定义的可行区域内运行；而经过足够长的时间窗口后，快速过程持续存在的残差
\[
\langle\varepsilon_f\rangle_T
\]
又会逐渐成为修改慢变量的证据。

这给出 AgentOS 中极为重要的双向规律：
\[
\boxed{
SlowStructureShapesFastActivity
}
\]
以及
\[
\boxed{
PersistentFastResidualShapesSlowStructure.
}
\]

例如，$\Theta$ 中已经形成的“门可以通过把手开启”这一结构，会直接改变工作记忆面对一扇门时的搜索空间，而不需要重新从视觉像素中探索所有可能动作。这是慢结构对快过程的约束。反过来，如果某类新式门在大量场景中持续使原有结构预测失败，那么情景残差会经过 $E$ 的积累逐渐推动 $\Theta$ 更新。这是快残差对慢结构的塑造。

在整个 AgentOS 中，我们可以得到一种并非严格固定、但具有结构意义的时间尺度关系：
\[
\tau_{\mathrm{servo}}
<
\tau_{\mathrm{BPCC}}
<
\tau_{\mathrm{SGC/FCS}}
<
\tau_h
<
\tau_E
<
\tau_\Theta
<
\tau_\Psi
<
\tau_\Omega.
\]
这里并不要求每一级都具有统一的绝对时间常数。例如，一个高危事件可能使 $\Theta$ 在短时间内发生较大更新，而一个长期重复任务也可能让某些 $E$ 长期保持可检索状态。真正需要保持的是相对时间尺度上的结构保护：越接近身份、价值和生命周期方向的变量，其修改证据应当越需要跨越更长时间窗口和更多独立经历。

这条原则可以称为\textbf{慢变量保护原则}：
\[
\boxed{
FastNoise
\nrightarrow
SlowGlobalStructure.
}
\]
否则，一次局部失败就可能改变长期身份，一个瞬时奖励就可能重写价值结构，一个短暂认知相态就可能被错误解释成长期人格。这会导致持续智能体无法形成真正意义上的结构连续性。

因此，时间层级并不是为了让 AgentOS 变得“像人脑”，而是一个稳定持续学习系统在数学上自然需要的条件：快速结构保证适应性，慢结构保证连续性，而不同时间尺度之间受控的信息交换则产生真正的长期学习。


\section{双轴状态空间：Agent-CSO 在功能与时间中的位置}

当功能层级和时间层级同时建立以后，我们就可以给 Agent-CSO 一个比“存在哪个模块中”更精确的描述。设 $C$ 为某一认知结构本体，Agent $A$ 在时刻 $t$ 对它的具体承载表示为
\[
AgentCSO_A(C,t)
=
(C,R_t,F_t,\tau_t,B_t,S_t),
\]
其中 $R_t$ 表示具体表示结构，$F_t$ 表示主要功能承载位置，$\tau_t$ 表示有效变化时间尺度，$B_t$ 表示内外边界及物理或数字载体，$S_t$ 表示激活、潜伏、预测、回忆等运行状态。

本章关注其中两个最基本坐标：
\[
(F_t,\tau_t).
\]
因此可以定义一个 Agent-CSO 双轴空间
\[
\mathcal X_{FT}
=
\mathcal F\times\mathcal T,
\]
其中 $\mathcal F$ 是功能空间，$\mathcal T$ 是时间尺度空间。智能体在任意时刻并不是简单拥有一组知识，而是具有某种 Agent-CSO 在 $\mathcal X_{FT}$ 中的分布：
\[
\rho(C,F,\tau,t).
\]
如果进一步对具体 CSO 做边缘化，可得到系统整体结构分布
\[
\rho_{CSO}(F,\tau,t).
\]

这个分布提供了一个观察成熟智能的新方式。例如，一个刚学习软件操作的 Agent，可能有大量任务结构集中在高显式认知区域：
\[
\rho_{CSO}
\quad
\text{集中于}
\quad
(h,\tau_h).
\]
它必须不断读界面、保持操作步骤、判断按钮意义并显式规划下一步。随着经验增加，同样的结构逐渐向稳定技能和局部控制区域迁移：
\[
(h,\tau_h)
\rightarrow
(\Theta,\tau_\Theta)
\rightarrow
(BodySkill,\tau_{\mathrm{fast}}).
\]
这里非常值得注意：最后一步同时包含两个方向不同的变化。认知结构先通过经验积累向更慢、更稳定的 $\Theta$ 沉降，然后由 $\Theta$ 编译出的技能却可以运行在比 $h$ 更快的局部闭环中。因此，“向慢学习”和“向低层执行”并不是同一个过程。

这正是双轴模型相对于单轴层级模型的优势。如果只画“高层—低层”，很容易认为越成熟的结构就越低层；如果只画“快—慢”，又很容易误认为越成熟的结构永远运行得越慢。真实情况是：一个技能的\textbf{结构形成}可能是慢过程，而它的\textbf{执行闭环}却可能是非常快的过程。

我们因此应该区分
\[
\tau_{\mathrm{formation}}
\]
和
\[
\tau_{\mathrm{execution}}.
\]
例如，一个熟练抓取技能可能经过数小时甚至数日经验形成，但一旦形成，其闭环控制可在毫秒或几十毫秒尺度执行。这意味着一个成熟 Agent-CSO 可以拥有慢速稳定的结构基础和快速局部的执行投影：
\[
C^{stable}_{\Theta}
\xrightarrow{\mathrm{compile}}
C^{procedural}_{Body}.
\]

由此，AgentOS 不再是一张静态模块架构图，而是一张不断发生结构迁移的双轴动力图。真正的长期学习，就是 $\rho_{CSO}(F,\tau,t)$ 随时间发生受控重新分布；真正的认知异常，则经常可以描述成某些结构长期停留在不合适的功能位置或时间尺度。例如，本应已经自动化的任务持续占据 $h$，表现为显式认知成本过高；本应重新显式检查的旧技能仍停留在局部闭环，则可能形成自动化错误；短时事件过快进入 $\Psi$，则可能导致身份漂移。

因此，本章的双轴理论实际上提供了一个统一的诊断问题：
\begin{statementbox}
这个 Agent-CSO 当前是否存在于正确的功能位置和时间尺度？
\end{statementbox}
对于长期智能而言，这个问题往往比“这个知识是否存在于模型中”更加重要。


\section{第一种跨尺度流：向慢巩固与结构沉降}

双轴架构中的第一类基本流，是认知结构从快速、易变、显式状态向慢速、稳定、可复用结构迁移。我们把它称为\textbf{向慢巩固（slowward consolidation）}。

在三相记忆中，它最直接表现为
\[
h\rightarrow E\rightarrow\Theta.
\]
当前工作记忆 $h$ 承载的是正在发生的 Agent-CSO。这些结构中只有部分具有长期价值，因此 $h\rightarrow E$ 并不是完整拷贝，而是选择性的经历写入。经过多个事件、多个上下文以及现实验证以后，$E$ 中反复出现的关系进一步被抽象和压缩为 $\Theta$ 中的生成结构：
\[
E
\xrightarrow{\mathrm{abstraction}}
\Theta.
\]

为了形式化这一过程，可以设某一 Agent-CSO $C$ 在时间窗口 $T$ 中出现的经验证据为
\[
\mathcal E_T(C)
=
\{e_1,\ldots,e_n\}.
\]
定义其结构稳定性估计
\[
\Sigma_T(C)
=
F
\left(
N_C,
Consistency_C,
Utility_C,
PredictiveGain_C,
Residual_C^{-1}
\right),
\]
其中 $N_C$ 表示重复出现次数，$Consistency_C$ 表示不同场景中的结构一致程度，$Utility_C$ 表示对目标闭合的贡献，$PredictiveGain_C$ 表示该结构对未来预测的提升，而 $Residual_C$ 表示使用该结构之后仍然产生的现实残差。

当
\[
\Sigma_T(C)>\theta_{\mathrm{consolidate}}
\]
并持续足够时间以后，该结构就具备向更慢层级沉降的条件。

这里需要特别强调：向慢巩固不是“把更多东西放进长期记忆”。真正重要的是结构自由度减少。一个经验集合可能包含大量细节：
\[
e_1,e_2,\ldots,e_n,
\]
而稳定结构则尝试提取
\[
\Theta_C
=
\arg\min_{\Theta}
\left[
L_{\mathrm{prediction}}
+
\lambda C_{\mathrm{representation}}
\right].
\]
也就是说，系统寻找一个既能解释过去经历又拥有更低表示负担的生成结构。从这个意义上讲，结构沉降同时是一种跨时间的压缩过程。

这种机制解释了为什么长期学习能够释放工作记忆。假设某一任务最初需要同时维持 $k$ 个显式关系：
\[
\{c_1,\ldots,c_k\}\subset h.
\]
经过结构化以后，它们可能被压缩为一个更高层生成结构
\[
C^*=\Phi(c_1,\ldots,c_k),
\]
未来只需要激活 $C^*$，就可以重新生成大量原有关系。于是原本的显式复杂性被慢结构吸收：
\[
C_h
\downarrow,
\qquad
C_\Theta
\uparrow.
\]

但按照本书后期确立的 CSO 本体，我们不把这种现象理解为某种“复杂度物质”真的从 $h$ 流入 $\Theta$。更加准确的是，同一批相关 Agent-CSO 改变了自己的表示粒度、稳定性和主要承载位置。所谓复杂度下降，是这种结构迁移在运行负载层面的结果。

向慢巩固因此具有两个相反但统一的效果：一方面，它使当前认知需要维护的自由度越来越少；另一方面，它使未来行为受到越来越多历史结构的约束。成熟 Agent 并不是“过去对现在影响越来越少”，恰恰相反，它的过去越来越深地进入未来，只是不再以显式回忆和逐步推理的形式出现，而是沉降为先验、技能、结构偏置与生成器。

因此可以概括为：
\[
\boxed{
PastExplicitProcessing
\rightarrow
PresentSlowStructure
\rightarrow
FutureFastConstraint.
}
\]


\section{第二种跨尺度流：向快激活与慢结构重入当前}

如果只有向慢巩固，智能体最终会成为越来越稳定但越来越难以适应现实变化的静态系统。因此 AgentOS 必须具有相反方向的第二种流：已经沉降到慢尺度中的结构，可以在当前需要时重新进入快速认知和行为过程。我们把它称为\textbf{向快激活（fastward activation）}。

最基础的形式包括
\[
E\rightarrow h
\]
和
\[
\Theta\rightarrow h.
\]
前者是情景召回：当前状态与过去某段经历形成结构相似性，于是过去轨迹的一部分重新进入当前认知。后者是结构激活：某种长期规律、技能、世界模型或预测结构直接对当前 $h$ 施加生成约束。

可以形式化定义当前状态
\[
h_t
\]
与结构记忆单元 $C_i^\Theta$ 的激活得分：
\[
A_i(t)
=
w_rR(C_i^\Theta,h_t)
+
w_gG(C_i^\Theta,G_t)
+
w_vV(C_i^\Theta,\Psi_t)
+
w_pP(C_i^\Theta,\varepsilon_t)
-
w_cCost(C_i^\Theta),
\]
其中 $R$ 表示当前状态相关性，$G$ 表示当前目标关联，$V$ 表示自我与价值相关性，$P$ 表示对当前残差的解释潜力，而 $Cost$ 表示将该结构提升进入显式认知的工作记忆成本。

只有当
\[
A_i(t)>\theta_{\mathrm{activate}}
\]
时，该结构才值得进入当前快速表面。

这一公式揭示了一个重要事实：成熟智能并不是“拥有越多长期结构，就同时激活越多结构”。恰恰相反，长期结构越丰富，选择性激活就越重要。因为 $h$ 的容量受到有限工作记忆定理约束：
\[
C_{\mathrm{run}}(t)\le C_{\max}(W).
\]
因此真正优秀的 AgentOS 必须同时拥有丰富的慢结构和稀疏的快速显式表面。

这里也解释了为什么召回与激活应当受到目标、自我和当前残差共同控制。一个结构虽然与当前输入存在语义相似性，却不一定与当前问题相关；一个过去经历虽然非常强烈，却不应因为记忆强度高就不断进入当前工作记忆。AgentOS 需要做的是从庞大的慢结构库中选择最有助于当前闭环闭合的少量结构。

向快激活也不仅意味着“回忆”。例如，$\Theta$ 中的熟练技能并不一定先进入文字化 $h$ 再指导动作。它可以直接投射到 Body Skill：
\[
\Theta
\rightarrow
BodySkill.
\]
此时发生的是慢结构向快速局部执行器的激活，而非慢结构向显式工作记忆的召回。这进一步说明，时间轴与功能轴必须分开建模。

如果一个局部技能运行顺利，则
\[
\Theta\rightarrow BodySkill
\]
已经足够；如果发生异常，则其相关结构可能进一步被提升：
\[
BodySkill
\rightarrow
Residual
\rightarrow
h.
\]
由此可以看到，向快激活不是单一通路，而是一组按照任务需要把稳定结构重新投射到不同快速功能层的机制。

这一机制最终让 Agent 获得一种重要性质：过去并不需要被不断显式记住，但过去能够在正确时刻重新参与现在。所谓长期记忆真正有价值的状态，并不是大量存储，而是
\[
\boxed{
SlowStructureCanReenterFastDynamicsWhenNeeded.
}
\]


\section{第三种跨尺度流：异常残差的功能向上升级}

前两种流主要沿时间尺度轴发生：当前结构向慢处巩固，慢结构又重新进入快速过程。然而智能体还面临另一个完全不同的问题：一个 Agent-CSO 当前由哪个功能层处理？如果局部功能可以解决问题，就没有必要交给更大的认知系统；只有当局部模型持续失败，或者事件具有更高层目标意义时，才应该向上升级。

因此，第三种基本流是\textbf{功能向上升级（functional upward promotion）}。

以身体控制为例，BPCC 对局部状态进行预测：
\[
\hat x_{t+\delta}
=
F_B(x_t,u_t).
\]
现实重新进入 Body Runtime 后得到
\[
x^{obs}_{t+\delta},
\]
于是形成局部残差
\[
\varepsilon_B
=
d
\left(
x^{obs}_{t+\delta},
\hat x_{t+\delta}
\right).
\]
如果
\[
\varepsilon_B<\theta_{\mathrm{local}},
\]
则这一扰动留在局部闭环中处理；若残差持续、风险升高或与高层目标相关，则形成升级指标
\[
\Gamma
=
F
\left(
Residual,
Persistence,
Novelty,
Risk,
GoalRelevance,
ModelUncertainty
\right).
\]
当
\[
\Gamma\ge\theta_{\mathrm{promote}},
\]
相关 Agent-CSO 被提升至更高功能层级。

因此：
\[
\boxed{
RealityDoesNotContinuouslyBecomeThought;
RelevantDeviationBecomesThought.
}
\]

这是 AgentOS 与简单“所有传感器数据都送进大模型”架构之间的根本区别。现实世界每一毫秒都在发生大量变化。如果所有变化都持续进入 $h$，那么工作记忆有限性将立即被突破。Body Runtime、工具、文件系统以及局部任务执行器存在的意义，就是让大量普通变化在局部形成闭合。只有那些无法被既有局部结构解释的残差，才具有成为全局认知事件的资格。

这一原则不仅适用于身体。在工具调用中，如果 API 正常返回预期结果，则无需重新规划；如果出现意外返回、权限错误或语义冲突，则需要提升。情景记忆召回同样如此：如果当前事件可以被已有结构稳定解释，就不需要重新审视 $\Theta$；如果重复经验不断产生结构性残差，才需要进行更慢的结构重解释。

因此向上升级的真正对象不是“数据”，而是\textbf{无法局部闭合的 Agent-CSO 及其残差关系}。可以写成：
\[
(C_i,\varepsilon_i)@F_k
\longrightarrow
(C_i,\varepsilon_i)@F_{k+1}.
\]

这种机制使 AgentOS 具备一种计算上极其重要的性质：显式认知是稀疏的。高层不负责持续监控所有底层细节，而只接收具有足够结构意义的偏差事件。于是注意力也可以重新定义：
\[
\boxed{
Attention
=
CrossScalePromotionPolicy.
}
\]
它并不只是选择当前文本 token 权重，而是决定哪些局部结构值得跨越功能层级进入稀缺的全局认知空间。

从控制理论来看，这类似事件触发控制；从操作系统角度看，它类似异常、中断和优先级提升；从认知角度看，它则表现为“某件原本自动进行的事情突然进入意识”。AgentOS 将这些现象统一到同一个原则之下：
\[
\boxed{
PredictLocally,\ EscalateResidualGlobally.
}
\]


\section{第四种跨尺度流：成熟能力的功能向下编译}

与异常向上相反，反复需要高层处理的结构如果长期稳定，就不应该永久占据昂贵的显式认知资源。第四种基本流因此是\textbf{功能向下编译（functional downward compilation）}。

假设某类任务 $T$ 最初需要工作记忆中的完整显式过程：
\[
h_t
=
\{
Goal,
State,
Plan,
Constraints,
ActionSequence
\}.
\]
在多次成功执行以后，情景记忆中会积累轨迹集合
\[
\mathcal E_T
=
\{
\tau_1,\tau_2,\ldots,\tau_n
\}.
\]
如果这些轨迹在状态条件和结果之间具有稳定映射，则系统可以从中提取策略结构
\[
\pi_T
=
\Phi(\mathcal E_T),
\]
并进一步形成
\[
BodySkill_T
=
(
Predictor_T,
Controller_T,
ResidualModel_T
).
\]
此时原本依赖显式规划的过程被编译为局部闭环：
\[
\boxed{
ExplicitCognitiveControl
\rightarrow
RepeatedTrajectory
\rightarrow
StablePolicy
\rightarrow
CompiledLocalLoop.
}
\]

因此向下编译不是简单“记住答案”。真正成熟的能力必须能够在现实反馈中自主关闭一个局部循环。如果某个 Skill 只是一串固定命令，而不能根据状态偏差调整执行，那么它仍然只是宏脚本，而不是完整意义上的程序化 Agent-CSO。

从 CSO 的角度看，向下编译可以描述为同一任务相关结构从高功能位置迁移至低功能位置：
\[
AgentCSO_T@h
\rightarrow
AgentCSO_T@F_{\mathrm{local}}.
\]
与此同时，结构形成的证据则沉降在较慢层：
\[
E_T
\rightarrow
\Theta_T.
\]
因此，一个成熟技能实际上具有“双重结构”：
\[
\Theta_T
\quad+\quad
F^{compiled}_T.
\]
前者是慢速生成和监督结构，后者是快速执行闭环。

这使我们能够严格解释“熟练以后为什么不用思考”。不是相关认知结构消失，而是其主要功能承载位置发生迁移：
\[
\boxed{
ThinkingLess
\neq
KnowingLess.
}
\]
更准确地说：
\[
\boxed{
StructureChangedCarrierAndFunctionalLocation.
}
\]

但功能向下编译必须保持可逆性。现实变化可能使过去稳定的技能失效，因此局部执行必须持续产生 residual。如果
\[
\varepsilon_{\mathrm{skill}}
>
\theta_{\mathrm{reopen}},
\]
就应触发
\[
F_{\mathrm{local}}
\rightarrow
h,
\]
重新进入显式认知。于是“能力向下”和“异常向上”形成一对完整的动态机制：
\[
\boxed{
Explicit
\rightleftarrows
Automatic.
}
\]

这一可逆性极其重要。没有向下编译，Agent 永远无法成熟；没有向上再显式化，Agent 又会被旧技能困住。真正的持续智能必须在自动化效率与重新认知能力之间保持动态平衡。


\section{Attention 作为跨尺度提升策略}

传统模型语境中的 attention 通常被理解为对当前表示空间中不同元素分配权重。在 AgentOS 中，这一概念必须扩展。因为持续智能面临的最稀缺资源并不是单纯计算量，而是\textbf{哪些结构能够进入有限全局工作区并获得跨系统协调权}。

因此，本章把 Attention 定义为：
\[
\boxed{
Attention
=
CrossScalePromotionPolicy.
}
\]

设系统当前存在候选 Agent-CSO 集合
\[
\mathcal C_t
=
\{C_1,\ldots,C_n\}.
\]
每个结构具有一个提升效用
\[
U_i
=
F
\left(
Relevance_i,
Novelty_i,
Residual_i,
Risk_i,
Goal_i,
Self_i,
Uncertainty_i
\right),
\]
同时占用工作记忆资源
\[
c_i.
\]
则 Attention 可以被形式化成有限容量下的选择问题：
\[
\max_{\mathcal S_t\subseteq\mathcal C_t}
\sum_{C_i\in\mathcal S_t}U_i
\]
满足
\[
\sum_{C_i\in\mathcal S_t}c_i
\le
C_{\max}(W).
\]

由此，注意不只是“聚焦哪个对象”，而是在决定：
\[
\boxed{
WhichLocalStructuresDeserveGlobalCognitiveAuthority?
}
\]

这解释了为什么 Boundary、SPC、CCM、Scheduler、$\Psi$ 都会参与 Attention 的形成。Boundary 决定外部结构是否具有进入资格；SPC 估计系统当前负载与相态；CCM 估计加入新结构以后是否突破工作记忆容量；Scheduler 决定不同目标和任务的优先级；$\Psi$ 则提供更慢的价值和身份相关性。因此 Attention 并不存在于单个模块中，而是多个尺度约束共同作用后形成的提升策略。

进一步说，Attention 还具有时间尺度选择功能。某件事情可能值得进入 $h$，但不值得写入 $E$；某段经历可能值得进入 $E$，但不值得修改 $\Theta$；某个结构变化可能值得修改 $\Theta$，却远不足以改变 $\Psi$。因此每一级慢化过程都有自己的“注意阈值”：
\[
\theta_h,
\quad
\theta_E,
\quad
\theta_\Theta,
\quad
\theta_\Psi.
\]
一般应满足越来越严格的慢变量保护要求：
\[
Evidence_{\Psi}
\gg
Evidence_{\Theta}
\gg
Evidence_E.
\]

这意味着 Attention 并不是只有当前时刻的注意力，还存在生命周期意义上的结构注意力：哪些事件值得成为经历，哪些经历值得成为能力，哪些能力值得成为自我定义的一部分。

从这一角度看，一个真正成熟的 Agent 并不是拥有无限注意，而是形成了一套越来越精确的跨尺度提升策略。它可以让绝大多数已知过程保持在局部，让真正重要的新颖结构进入工作记忆，让持续重复的经验进入慢结构，同时保护最慢的身份和生成先验不被瞬时噪声扰动。

因此注意力真正解决的是：
\[
\boxed{
FiniteGlobalCapacity
\quad\text{vs.}\quad
InfinitePotentialWorldEvents.
}
\]
它是有限主体能够在开放世界中持续存在的核心选择机制。


\section{成熟智能：选择性显式认知而非持续中央思考}

双轴架构最终给出了一个与“大模型越强就应该思考得越多”完全不同的智能成熟观。若一个 Agent 每面对任何熟悉问题仍然必须重新加载大量上下文、执行完整链式推理、调用中央模型并显式维护所有细节，那么它即使单次推理能力很强，仍然是一种结构上不成熟的智能。

成熟智能真正应该出现的趋势是：
\[
\boxed{
ExplicitGlobalInterventionRatio
\downarrow
}
\]
同时：
\[
\boxed{
SuccessfulLocalClosureRatio
\uparrow.
}
\]

可以定义一个概念性的显式认知率
\[
R_{\mathrm{explicit}}
=
\frac{
T_{\mathrm{global\ explicit}}
}{
T_{\mathrm{successful\ behavior}}
},
\]
或者按事件数定义
\[
R_{\mathrm{explicit}}
=
\frac{
N_{\mathrm{events\ entering}\ h}
}{
N_{\mathrm{behaviorally\ relevant\ events}}
}.
\]
当一个领域逐渐熟练时，在行为可靠性不下降的情况下，$R_{\mathrm{explicit}}$ 应当下降。

这与人的熟练行为具有明显的结构相似性：初学驾驶时，大量视觉对象、规则、动作顺序都进入显式工作记忆；熟练以后，大量低层操作在局部闭环中完成，而显式注意主要处理路线选择、异常车辆、复杂交通和新环境。但本理论并不依赖人类认知类比。任何受有限全局容量约束的长期智能，只要需要在高频环境中保持高效运行，就会面临类似的层级化压力。

因此：
\[
\boxed{
MatureIntelligence
=
SelectiveExplicitCognition.
}
\]

所谓“选择性”，至少包含三层含义。第一，正常局部动力不进入 $h$；第二，稳定经验不长期占据 $h$，而应沉降到 $E/\Theta$ 或局部 Skill；第三，即使长期记忆非常丰富，也只有与当前目标、风险和残差有关的部分才重新进入显式认知。

这种结构还解释了为什么有限工作记忆不是智能的缺陷，而可能是高级智能组织化的动力来源。如果工作记忆无限，系统没有足够压力去形成抽象、技能、局部闭环、工具和外部记忆；正因为
\[
C_h<C_{\mathrm{world}},
\]
系统才必须学习
\[
Compression,
\quad
Abstraction,
\quad
Automation,
\quad
Offloading,
\quad
Hierarchy.
\]
因此有限性促成结构。

当然，显式认知率下降不能被误解为“越少思考越聪明”。如果一个系统在新环境中仍然拒绝提升异常结构，那么它只是僵化。成熟的关键不是显式认知少，而是：
\[
\boxed{
ExplicitWhenNecessary,
LocalWhenSufficient.
}
\]

因此我们可以把成熟度写成一个更完整的函数：
\[
M
=
F
\left(
LocalClosure,
PromotionPrecision,
CompilationQuality,
RecoveryAbility,
RealityAlignment
\right).
\]
其中任何一个维度缺失都可能产生伪成熟：局部闭合率高但现实错误率也高，是自动化僵化；升级率很高但局部闭合率低，是中央认知过载；技能编译很好但无法重新显式化，则缺乏适应性。

成熟智能不是“思考最多的系统”，而是\textbf{知道什么不需要思考、什么必须重新思考，以及什么经验值得永久改变未来思考方式的系统}。


\section{学习作为双轴空间中的结构重新分布}

在前述理论中，我们已经逐渐放弃把学习简单定义成参数改变。参数改变可以是学习的一个具体实现，但从 AgentOS 的系统层看，更一般的学习应当描述为 Agent-CSO 在表示、功能位置、时间尺度和载体之间发生受控迁移。

对于本章的双轴模型，我们可以把这种学习首先写成：
\[
(F_t,\tau_t)
\rightarrow
(F_{t+\Delta t},\tau_{t+\Delta t}).
\]
若考虑全部 Agent-CSO，则系统学习对应整体分布变化：
\[
\rho_{CSO}(F,\tau,t)
\rightarrow
\rho_{CSO}(F,\tau,t+\Delta t).
\]

于是可以定义一个双轴结构流
\[
\mathbf J_{CSO}
=
\left(
J_F,
J_\tau
\right),
\]
其中 $J_F$ 表示功能轴迁移流，$J_\tau$ 表示时间尺度轴迁移流。理论上可以写一个结构连续方程：
\[
\frac{\partial \rho_{CSO}}{\partial t}
+
\nabla_F\cdot J_F
+
\frac{\partial J_\tau}{\partial \tau}
=
S_{CSO}
-
D_{CSO},
\]
其中 $S_{CSO}$ 表示新认知结构形成，$D_{CSO}$ 表示遗忘、合并、压缩和结构消解。

必须再次强调，这里并不存在严格的物理“CSO 守恒”。一个结构可以与其他结构合并，一个抽象结构可以代表大量具体结构，一个错误结构也可以被彻底废弃。因此该方程的意义是提供一种结构流描述语言，而不是宣称认知结构具有类似质量或能量的守恒定律。

有了这一框架，我们可以把多种看似不同的学习现象统一起来：

\[
h\rightarrow E
\]
是快速状态向较慢轨迹结构迁移；

\[
E\rightarrow\Theta
\]
是经验向生成结构迁移；

\[
\Theta\rightarrow h
\]
是慢结构向快速显式认知重新激活；

\[
h\rightarrow Skill
\]
是高层显式结构向局部执行功能迁移；

\[
Skill\rightarrow h
\]
则是技能异常后的重新显式化；

\[
Internal\rightarrow Tool
\]
是结构跨边界迁移；

而长期稳定的
\[
F_i\rightarrow F_j
\]
迁移如果持续出现，则可能进一步沉降为新的功能连接。

因此：
\[
\boxed{
Learning
=
StructuralRedistributionAcrossScales.
}
\]

这一定义比“模型权重更新”更适合持续 Agent。因为一个 Agent 完全可以在基础模型参数几乎不变的情况下，通过写入 $E$、重构 $\Theta$、形成 Skill、建立工具、改变 Boundary 策略和重新分配功能承载位置而发生显著学习。反过来，即使基础模型参数被持续更新，如果所有任务仍然需要同一种中央推理过程，也不意味着系统真正发生了成熟意义上的结构发展。

这一观点同时把学习与发育区分开来。普通学习改变
\[
\rho_{CSO}(F,\tau,t),
\]
而更长期的发育还可能改变功能拓扑本身：
\[
\mathcal G_F(t)
\rightarrow
\mathcal G_F(t+\Delta t).
\]
在当前固定总体拓扑的 AgentOS 中，我们首先依靠已有功能节点内部的承载迁移实现大部分持续学习；只有当长期结构失配无法再通过现有迁移解决时，才需要考虑更深层拓扑重构。因而双轴迁移理论既解释了固定拓扑 AgentOS 能够获得相当强的持续适应能力，也为更长期的结构发育留下了清晰接口。


\section{双轴架构的稳定性原则与统一动力学}

到这里，我们已经得到四种基本跨尺度流：
\[
\boxed{
\begin{aligned}
&\text{时间向慢：} &&Consolidation,\\
&\text{时间向快：} &&Activation,\\
&\text{功能向上：} &&ResidualPromotion,\\
&\text{功能向下：} &&Compilation.
\end{aligned}
}
\]
真正的 AgentOS 不只是允许这些流存在，还必须保证它们不会形成跨尺度振荡和结构污染。因此双轴架构需要一组统一稳定性原则。

第一是\textbf{局部闭合原则}：
\[
\boxed{
NormalDynamicsStayLocal.
}
\]
任何已经可以在较低功能层级可靠解决的过程，不应持续占用高层工作记忆。

第二是\textbf{慢约束原则}：
\[
\boxed{
SlowGovernance
\neq
FastMicromanagement.
}
\]
$\Theta$、$\Psi$ 和 $\Omega$ 应通过边界、先验、目标偏置和可行域约束快速系统，而不是逐步决定每一个局部动作。否则慢变量不仅失去稳定性，还会成为整个系统的实时瓶颈。

第三是\textbf{残差提升原则}：
\[
\boxed{
PersistentImportantResidual
\Rightarrow
Promotion.
}
\]
局部偶发噪声不应提升，但持续、新颖、高风险或与高层目标相关的残差必须具有进入更高认知层的路径。

第四是\textbf{结构沉降原则}：
\[
\boxed{
RepeatedStableCSORelation
\Rightarrow
MoreStableOrLocalCarrier.
}
\]
如果某种认知结构长期以同一种关系重复被显式处理，那么系统应尝试压缩、结构化、编译或外化，而不是永久维持高层重复计算。

第五是\textbf{慢变量保护原则}：
\[
\boxed{
FastNoise
\nrightarrow
SlowGlobalStructure.
}
\]
不同层级之间必须存在证据累积时间窗和迟滞，否则一次局部异常就可能层层放大成身份和长期方向变化。

第六是\textbf{现实闭合原则}：
\[
\boxed{
Reality
=
UltimateConstraint.
}
\]
无论一个 Agent-CSO 在内部经历多少层激活、压缩、预测和重新解释，其关于外部现实的有效性最终都必须通过行动—观测闭环重新接受现实约束。预测不能冒充观测，内部一致性不能代替世界验证，自我模型也不能成为不可证伪的封闭系统。

在数学上，我们可以把完整系统抽象成：
\[
\dot x
=
F(x,z,\mathcal G_F,u,w),
\]
其中 $x$ 表示快速状态；$z$ 表示慢结构；$\mathcal G_F$ 表示功能拓扑；$u$ 表示来自 TCE、Scheduler、Boundary 等治理结构的控制；$w$ 表示外部世界扰动。

慢结构满足：
\[
\dot z
=
\epsilon
G
\left(
z,
\langle \varepsilon(x,w)\rangle_T,
E,
\Psi
\right),
\]
而 Agent-CSO 承载分布满足：
\[
\frac{\partial\rho}{\partial t}
=
-\nabla\cdot\mathbf J
+
S-D.
\]

因此 AgentOS 的长期轨迹并不是所有局部轨迹的简单线性和：
\[
\Gamma_{\mathrm{Agent}}
\neq
\sum_i\Gamma_i.
\]
更加准确地说：
\[
\boxed{
\Gamma_{\mathrm{Agent}}
=
\Pi_{\mathrm{behavior}}
\left[
\mathcal D_{\mathrm{multiscale}}
\bowtie
Environment
\right],
}
\]
即我们最终观察到的主体行为，是多时间尺度、多功能闭环和现实环境长期非线性耦合后的投影。

这也是为什么不能通过单独分析 LLM、Memory、Body 或 Scheduler 中任意一个组件来解释整个 Agent。持续智能的关键性质正产生于这些闭环之间的相对速度、局部闭合、结构沉降和受控升级。


\section{统一结论：AgentOS 是功能层级与时间层级的乘积空间}

现在我们可以重新回到本章开头的基本命题：
\[
\boxed{
AgentOS
=
FunctionalHierarchy
\times
TemporalHierarchy.
}
\]
这一公式并不是为了提供一个漂亮的架构口号，而是指出 AgentOS 中任何关键结构都不能只用“模块归属”描述。一个 Agent-CSO 真正的系统位置至少必须同时包含：
\[
\boxed{
WhereItIsProcessed
+
HowFastItChanges.
}
\]

功能层级决定它由哪个闭环承担；时间层级决定它在多长时间范围内保持稳定。学习则改变这种位置：
\[
(F,\tau)
\rightarrow
(F',\tau').
\]

由此，我们得到整个 AgentOS 最基础的四向动力学：

\[
\boxed{
\begin{array}{ccc}
& \text{功能向上：异常/新颖结构升级} & \uparrow\\[4pt]
\text{时间向慢：巩固}
& AgentCSO &
\text{时间向快：激活}\\[4pt]
& \text{功能向下：成熟能力编译} & \downarrow
\end{array}
}
\]

这四种流共同解释了一个持续 Agent 为什么既可以成长又保持稳定。

如果没有向慢巩固，经验不能成为结构；
如果没有向快激活，过去不能重新参与现在；
如果没有功能向上升级，局部错误不能重新进入认知；
如果没有功能向下编译，成熟能力永远占据昂贵的显式处理资源。

因此四者缺一不可：
\[
\boxed{
PersistentIntelligence
=
Consolidation
+
Activation
+
Promotion
+
Compilation.
}
\]

但这四种动力又受到两个更基础的约束。

第一个来自主体内部：
\[
\boxed{
FiniteExplicitCapacity.
}
\]
工作记忆 $h$ 的有限性迫使系统不断进行选择、压缩、沉降和编译。

第二个来自主体外部：
\[
\boxed{
RealityConstraint.
}
\]
现实不断通过行动结果和预测残差检验内部结构是否仍然有效。

正是在这两个约束之间，AgentOS 获得了一种不同于传统静态软件和单次模型推理的生命周期动力学：有限工作记忆使它不能同时显式承担所有结构；慢记忆使过去能够保存；局部功能闭环使成熟能力不再重复占用全局认知；异常提升使系统能够重新关注失效结构；而现实反馈则持续决定哪些结构应继续存在、改变位置或重新学习。

由此我们可以给出本章最终命题：

\begin{statementbox}
成熟智能不是一个不断扩大的中央推理器，而是一个能够让认知结构在功能层级与时间层级之间持续找到适当承载位置，并能够在稳定、激活、异常与自动化之间可逆迁移的多尺度动力系统。
\end{statementbox}

进一步，如果我们把学习理解为 Agent-CSO 的双轴重新分布：
\[
\boxed{
Learning
=
RedistributionOfAgentCSO
AcrossFunctionAndTimescale,
}
\]
那么所谓“能力增长”就不再只是内部知识数量增加，而表现为越来越多过去必须昂贵显式处理的问题，被转化为稳定结构、局部技能、外部工具和环境结构；与此同时，系统仍然保留把失效结构重新提升到显式认知层进行修正的能力。

因此，一个真正成熟的 AgentOS 应呈现出一种看似矛盾但实际上高度统一的特征：
\begin{statementbox}
内部结构越来越丰富，显式认知却越来越稀疏。
\end{statementbox}
丰富的是跨尺度沉降下来的长期结构，稀疏的是每一个瞬间真正需要全局协调的 Agent-CSO。前者不断扩大智能体能够自然完成的行为范围，后者保护有限工作记忆不被整个世界的结构复杂性淹没。

最终，我们得到一个贯穿本章、也能够作为 AgentOS 多尺度架构基本定律的表达：

\begin{statementbox}
StableStructuresMoveTowardStableCarriers;
UnexpectedResidualsMoveTowardGlobalCognition;
SlowStructuresConstrainFastDynamics;
PersistentFastResidualsRewriteSlowStructures.
\end{statementbox}

即：

\begin{statementbox}
稳定结构向稳定承载处沉降，异常结构向全局认知层提升；
慢结构约束快速动力，持续快速残差反过来塑造慢结构。
\end{statementbox}

这四句话共同构成了 AgentOS 双轴多尺度架构的动力学核心。它使我们不再从“有多少模块”理解智能，而开始从“结构如何跨尺度存在、迁移、沉降、激活与重新进入现实”理解一个持续智能体。